\begin{textsample}{POS Dim 1 – vem\_ed\_30 – Score 24.00 – vem\_ed\_30/t158.txt}  \label{ex:f1_pos_027}
PCI aborda emissões de poluentes em veículos No primeiro semestre de 2025, o professor Marco Aurélio Fróes, do curso de Mecânica Automobilística da Fatec Santo André, desenvolveu um Projeto Colaborativo Internacional (PCI) com o professor Javier Pavón, da DUOC UC (Chile).
O projeto \textbf{permitiu} aos estudantes brasileiros e chilenos conhecerem normas internacionais \textbf{utilizadas} na certificação de emissões de poluentes em veículos leves, praticarem ensaios em laboratório e \textbf{compararem} os resultados.
O professor Orlando de Salvo Júnior, especialista em emissões veiculares, \textbf{foi} convidado pelo professor Fróes para auxiliar no entendimento dos \textbf{conteúdos} técnicos.
Esse \textbf{foi} o primeiro PCI de Fróes. “Fiquei um pouco perdido inicialmente, pois me faltava experiência, mas a professora Regiane orientou e conduziu muito bem, trazendo segurança”, comenta, referindo-se a Regiane Souza Camargo Moreira, responsável pela \textbf{coordenação} dos PCIs em espanhol na área de Apoio à Internacionalização do Ensino Superior da Divisão de Extensão e Pesquisa (Depes) da Coordenadoria Geral de Ensino Superior de Graduação (CGESG) do CPS.
As \textbf{interações} dos grupos de alunos se deram por WhatsApp, videochamadas no Zoom ou Google Meet, e os estudos desenvolvidos \textbf{foram} \textbf{compartilhados} no Padlet (em português e espanhol para driblar as barreiras idiomáticas).
No entender de Fróes, “a \textbf{troca} de experiências e de \textbf{realidades} com uma instituição de ensino internacional \textbf{foi} o principal ganho. \textbf{Foi} \textbf{possível} ver, na prática, a aplicação de uma norma internacional para homologação veicular com \textbf{todos} os critérios de segurança e condicionamento, bem como \textbf{estabelecer} comparações com os procedimentos de ensaio de potência e torque”. \textbf{continuação} O professor Javier Pavón, em sua \textbf{reflexão} final sobre o projeto \textbf{publicada} no Padlet, avaliou: “Durante \textbf{todo} o \textbf{processo}, os alunos adquiriram novos \textbf{conhecimentos} sobre normas e procedimentos relacionados a testes de motores em dinamômetro, o \textbf{que} lhes \textbf{permitiu} \textbf{ampliar} seus \textbf{conhecimentos} técnicos e aplicar metodologias de teste de \textbf{acordo} com as normas internacionais.
Além disso, o trabalho em conjunto com os colegas brasileiros ofereceu a possibilidade de conhecer diferentes culturas e experiências no setor automotivo, \textbf{promovendo} uma visão global sobre a homologação e certificação de veículos”.

% matched lemmas: acordo, ampliar, comparar, compartilhar, conhecimento, conteúdo, continuação, coordenação, estabelecer, interação, permitir, possível, processo, promover, publicar, que, realidade, reflexão, ser, todo, troca, utilizar
\end{textsample}
